\section{The local-time condition by construction}
\label{sec:local-time-by-construction}

This section proves that the deep ReLU input-convex neural network used in the
certificate architecture enforces the required sign of the local-time term by
construction.  The relevant point is a sign convention: the network
\(C_\psi\) defined below is convex, and hence its gradient jump across an
activation interface has the convex sign.  Therefore the supermartingale
certificate uses the residual form
\[
    V_\theta(x)=S_\theta(x)-C_\psi(x),
\]
where \(S_\theta\) is continuously differentiable.  The branch
\(-C_\psi\), rather than \(C_\psi\) itself, then satisfies the nonpositive
normal-derivative jump condition required by the It\^o--Tanaka formula.

\subsection{The deep ReLU ICNN}

Let \(x\in\mathbb R^d\).  Consider a network with \(L\geq 1\) hidden layers,
all of width \(m\), defined recursively by
\begin{align}
    z_1(x)
    &=
    \operatorname{ReLU}(U_1x+b_1),
    \\
    z_{\ell}(x)
    &=
    \operatorname{ReLU}
    \bigl(U_{\ell}x+b_{\ell}+W_{\ell}z_{\ell-1}(x)\bigr),
    \qquad \ell=2,\dots,L,
\end{align}
and with scalar output
\begin{equation}
    C_\psi(x)
    =
    u^\top x+c+w^\top z_L(x).
    \label{eq:deep-relu-icnn}
\end{equation}
The input matrices \(U_\ell\), biases \(b_\ell\), output input weight \(u\),
and output bias \(c\) are unrestricted.  The recurrent matrices and final
hidden-to-output weights satisfy
\begin{equation}
    W_\ell\geq 0,
    \qquad
    w\geq 0,
    \label{eq:icnn-nonnegative-weights}
\end{equation}
entrywise.

In the implementation, these inequalities hold for every parameter value
because
\[
    W_\ell=\operatorname{softplus}(\widetilde W_\ell),
    \qquad
    w=\operatorname{softplus}(\widetilde w).
\]
Since \(\operatorname{softplus}(r)>0\) for every finite \(r\), the effective
weights are nonnegative without requiring a constrained optimization step.

\begin{proposition}[Convexity and piecewise-affine structure]
\label{prop:deep-relu-icnn-convex}
The function \(C_\psi:\mathbb R^d\to\mathbb R\) defined by
\eqref{eq:deep-relu-icnn}--\eqref{eq:icnn-nonnegative-weights} is finite,
continuous, convex, and piecewise affine.
\end{proposition}

\begin{proof}
We prove convexity by induction over the hidden layers.

For the first layer, each coordinate has the form
\[
    (z_1)_k(x)
    =
    \max\{0,\,a_k^\top x+\beta_k\},
\]
which is the maximum of two affine functions and is therefore convex.
Consequently every coordinate of \(z_1\) is convex.

Suppose that every coordinate of \(z_{\ell-1}\) is convex.  For a fixed
coordinate \(k\) of the next preactivation, write
\[
    q_{\ell,k}(x)
    =
    (U_\ell x+b_\ell)_k
    +
    \sum_{r=1}^m (W_\ell)_{kr}(z_{\ell-1})_r(x).
\]
The first term is affine.  Every coefficient \((W_\ell)_{kr}\) is
nonnegative, so the second term is a nonnegative linear combination of convex
functions.  Thus \(q_{\ell,k}\) is convex.  Because ReLU is convex and
nondecreasing,
\[
    (z_\ell)_k(x)
    =
    \operatorname{ReLU}(q_{\ell,k}(x))
\]
is convex.  This completes the induction.

Finally,
\[
    C_\psi(x)
    =
    u^\top x+c+
    \sum_{k=1}^m w_k(z_L)_k(x)
\]
is the sum of an affine function and a nonnegative linear combination of
convex functions, and is therefore convex.

Every operation in the network is obtained from affine maps, finite sums, and
coordinatewise ReLU.  Hence the network is continuous and piecewise affine,
with finitely many affine pieces.
\end{proof}

\subsection{Convexity determines the interface sign}

Let \(K_i\) and \(K_j\) be two full-dimensional affine regions of
\(C_\psi\) that share a codimension-one face \(F\).  On the two regions,
write
\[
    C_\psi(x)=a_i^\top x+c_i,
    \qquad x\in K_i,
\]
and
\[
    C_\psi(x)=a_j^\top x+c_j,
    \qquad x\in K_j.
\]
Let \(n_{i\to j,F}\) be a unit normal to \(F\) pointing from \(K_i\) into
\(K_j\).

\begin{lemma}[Normal derivative jump of a convex CPWL function]
\label{lem:convex-normal-jump}
For every shared face \(F\),
\begin{equation}
    (a_j-a_i)\cdot n_{i\to j,F}\geq 0.
    \label{eq:convex-normal-jump}
\end{equation}
Equivalently,
\[
    \bigl(\nabla C_{\psi,j}(x)-\nabla C_{\psi,i}(x)\bigr)
    \cdot n_{i\to j,F}
    \geq 0,
    \qquad x\in\operatorname{relint}(F).
\]
\end{lemma}

\begin{proof}
Choose \(x_0\in\operatorname{relint}(F)\).  Since the two affine pieces agree
on \(F\), their difference vanishes on every tangent direction to \(F\).
Therefore \(a_j-a_i\) is normal to \(F\), so there exists
\(\lambda\in\mathbb R\) such that
\[
    a_j-a_i=\lambda n_{i\to j,F}.
\]
It remains to prove \(\lambda\geq 0\).

For sufficiently small \(\varepsilon>0\), the points
\[
    x_-=x_0-\varepsilon n_{i\to j,F}\in K_i,
    \qquad
    x_+=x_0+\varepsilon n_{i\to j,F}\in K_j.
\]
Restrict \(C_\psi\) to the line
\[
    g(t)=C_\psi(x_0+t n_{i\to j,F}).
\]
Because \(C_\psi\) is convex, \(g\) is a convex function of one variable.
Its left and right derivatives at zero are
\[
    g'_-(0)=a_i\cdot n_{i\to j,F},
    \qquad
    g'_+(0)=a_j\cdot n_{i\to j,F}.
\]
The one-dimensional derivative of a convex function is nondecreasing, hence
\[
    a_i\cdot n_{i\to j,F}
    \leq
    a_j\cdot n_{i\to j,F}.
\]
This is exactly \eqref{eq:convex-normal-jump}.
\end{proof}

\subsection{The residual certificate}

We now consider the certificate
\begin{equation}
    V_\theta(x)
    =
    \rho\bigl(S_\theta(x)-C_\psi(x)\bigr),
    \qquad \rho>0,
    \label{eq:residual-icnn-certificate}
\end{equation}
where \(S_\theta\in C^1\).  In the implementation, \(S_\theta\) is the
shallow piecewise-quadratic network
\[
    S_\theta(x)
    =
    c+p^\top x+\frac12x^\top Hx
    +\frac12\sum_{k=1}^M
    v_k\operatorname{ReLU}(a_k^\top x+b_k)^2.
\]
Since the scalar function \(r\mapsto \tfrac12\operatorname{ReLU}(r)^2\) is
continuously differentiable, \(S_\theta\) has no gradient jump across any of
its activation hyperplanes.

The common refinement of the activation regions of \(S_\theta\) and
\(C_\psi\) is a finite polyhedral partition.  Let \(R_i\) and \(R_j\) be
adjacent full-dimensional cells of this refinement, sharing a face \(F\), and
let \(n_{i\to j,F}\) point from \(R_i\) to \(R_j\).

\begin{theorem}[Local-time condition by construction]
\label{thm:deep-relu-icnn-local-time}
For every choice of the unconstrained parameters of \(S_\theta\) and
\(C_\psi\), every choice of raw recurrent and output parameters, and every
adjacent pair of cells in the common refinement, the certificate
\eqref{eq:residual-icnn-certificate} satisfies
\begin{equation}
    \bigl(
        \nabla V_{\theta,j}(x)-\nabla V_{\theta,i}(x)
    \bigr)
    \cdot n_{i\to j,F}
    \leq 0,
    \qquad x\in\operatorname{relint}(F).
    \label{eq:local-time-condition-by-construction}
\end{equation}
Consequently, every codimension-one local-time contribution in the
multidimensional It\^o--Tanaka formula for \(V_\theta(X_t)\) is nonpositive.
No activation-region enumeration or pairwise interface verification is needed
to establish this sign.
\end{theorem}

\begin{proof}
On the two sides of \(F\),
\[
    \nabla V_{\theta,i}
    =
    \rho\bigl(\nabla S_{\theta,i}-\nabla C_{\psi,i}\bigr),
    \qquad
    \nabla V_{\theta,j}
    =
    \rho\bigl(\nabla S_{\theta,j}-\nabla C_{\psi,j}\bigr).
\]
Therefore
\begin{align}
    &\bigl(
        \nabla V_{\theta,j}-\nabla V_{\theta,i}
    \bigr)\cdot n_{i\to j,F}
    \\
    &\qquad=
    \rho
    \bigl(
        \nabla S_{\theta,j}-\nabla S_{\theta,i}
    \bigr)\cdot n_{i\to j,F}
    -
    \rho
    \bigl(
        \nabla C_{\psi,j}-\nabla C_{\psi,i}
    \bigr)\cdot n_{i\to j,F}.
    \label{eq:residual-jump-decomposition}
\end{align}

Because \(S_\theta\in C^1\), its gradient is continuous across every face of
the common refinement.  Hence
\[
    \nabla S_{\theta,j}(x)-\nabla S_{\theta,i}(x)=0,
    \qquad x\in F.
\]
If \(F\) is introduced only by the smooth branch, then
\(C_\psi\) is affine on both sides and its gradient jump is also zero.  If
\(F\) lies on an activation interface of \(C_\psi\), Lemma
\ref{lem:convex-normal-jump} gives
\[
    \bigl(
        \nabla C_{\psi,j}(x)-\nabla C_{\psi,i}(x)
    \bigr)\cdot n_{i\to j,F}
    \geq 0.
\]
Substituting these facts into \eqref{eq:residual-jump-decomposition}, and using
\(\rho>0\), yields
\[
    \bigl(
        \nabla V_{\theta,j}(x)-\nabla V_{\theta,i}(x)
    \bigr)\cdot n_{i\to j,F}
    =
    -\rho
    \bigl(
        \nabla C_{\psi,j}(x)-\nabla C_{\psi,i}(x)
    \bigr)\cdot n_{i\to j,F}
    \leq 0.
\]
This proves \eqref{eq:local-time-condition-by-construction}.

In the multidimensional It\^o--Tanaka formula, the contribution of \(F\) is
\[
    \frac12
    \int_0^t
    \bigl(
        \nabla V_{\theta,j}(X_s)-\nabla V_{\theta,i}(X_s)
    \bigr)
    \cdot n_{i\to j,F}\,dL_s^F.
\]
Surface local time is nondecreasing, so \(dL_s^F\) is a nonnegative measure.
The integrand is nonpositive by the preceding argument; therefore the entire
integral is nonpositive.  Summing over all faces preserves this sign.
\end{proof}

\begin{corollary}[Max-affine specialization]
\label{cor:max-affine-local-time}
The conclusion of Theorem~\ref{thm:deep-relu-icnn-local-time} also holds when
the convex branch is
\[
    C_\psi(x)=\max_{1\leq r\leq R}\{a_r^\top x+b_r\}.
\]
In particular, it applies to the residual max-affine certificate used in the
initial radial Ornstein--Uhlenbeck training experiment.
\end{corollary}

\begin{proof}
A finite maximum of affine functions is finite, continuous, convex, and
piecewise affine.  The proof of
Theorem~\ref{thm:deep-relu-icnn-local-time} uses no other property of
\(C_\psi\), so it applies unchanged.
\end{proof}

\begin{remark}[What is and is not guaranteed]
The theorem concerns only the singular, interface part of the generalized
It\^o formula.  It does not by itself imply that \(V_\theta(X_t)\) is a
supermartingale.  One must still verify the interior generator inequality,
for example
\[
    \mathcal LV_{\theta,i}(x)\leq 0
\]
on every relevant smooth cell, together with the required boundary,
separation, stopping, and stochastic-integrability conditions.  The
architecture removes the combinatorial verification of the local-time sign;
it does not remove the ordinary generator verification problem.
\end{remark}

\begin{remark}[Distributional-Hessian interpretation]
The result can be stated without referring to a particular activation-region
partition.  A finite convex piecewise-affine function has a positive
semidefinite singular Hessian measure.  Since \(S_\theta\in C^1\), it has no
codimension-one singular Hessian mass, and therefore
\[
    D^2_{\mathrm{sing}}V_\theta
    =
    -\rho D^2C_\psi
    \preceq 0.
\]
Thus the architecture enforces nonpositive singular curvature globally.  The
facewise normal-jump condition in
\eqref{eq:local-time-condition-by-construction} is the polyhedral expression
of this measure inequality.
\end{remark}
